\section{The tiered access path: from collection to admitted segment}
\label{sec:retrieval}
\label{sec:lanes}

The access path reads the collections, catalogue and indices maintained by
acquisition and organisation. It is the query-side counterpart of the
management levels: the catalogue tier chooses a collection, the shelf tier
chooses documents within it, the document tier opens a body, and the segment
tier scores its parts. Three access policies use different affordances of that
path, and a calling workflow may invoke the separate server-side Laya evidence
gate. The tool fetches referenced document text, segments it and applies a
local, query-conditioned verdict. Its model-card field (\texttt{noul}) is a
typed 0--1 score; this study reports no probability-calibration experiment and
does not interpret the score as a calibrated probability. Segments meeting the
configured threshold (0.5 by default) are returned as survivors with
provenance. A scored segment below the threshold is not admitted, while tool
or fetch failures are reported separately; a partially completed call may
still return successfully scored survivors. A scoped not-found report is
reserved for a completed search with no supporting evidence, not for an
execution failure.

An access policy runs the same five steps and differs in the tiers it visits
and the tools it calls.
\begin{enumerate}
\item \emph{Rewrite.} The agent decomposes the question into a subject, an
attribute and any constraint (Lane A), or classifies it as literal or
enumerative for the scanner and semantic for the judge (Lane B).
\item \emph{Gather.} The policy collects candidates: dense top-30 above cosine
0.35 inside a pinned collection (A); catalogue pages followed by the scanner
or the judge (B); one server call that merges vector top-30 with catalogue
results (C).
\item \emph{Merge and cap.} Candidates are deduplicated by collection and
document, ordered with documents returned by both sources first, and capped at
twelve candidates.
\item \emph{Judge.} The gate fetches each candidate's text, splits it into
segments and scores every segment against the posed evidence question. Scored
survivors carry provenance; a fetch or tool failure is reported separately from
a low score.
\item \emph{Cut and read.} The answer step reads the first five survivors in
full. The shipped configuration orders them by verdict score; the filter replay
of Section~\ref{sec:e4} keeps the fusion order instead, and the relative margin
is evaluated there as a diagnostic rather than as a configured control.
\end{enumerate}

\begin{algorithm}[t]
\caption{Tiered access with optional evidence gate}
\label{alg:tiered}
\begin{algorithmic}[1]
\Require question $q$; catalogue $\mathcal{C}$; threshold $\tau$; gate flag $g$
\State $c \gets \textsc{SelectCollection}(\mathcal{C}, q)$ \Comment{catalogue tier}
\State $D \gets \textsc{Shelf}(c, q)$ \Comment{document descriptions}
\State $B \gets \textsc{OpenBody}(D)$ \Comment{document tier}
\State $S \gets \textsc{Segments}(B)$; score $S$ (dense + lexical) \Comment{segment tier}
\If{$g$}
  \State $S \gets \textsc{LayaJudge}(S, q, \tau)$; keep survivors with provenance
\EndIf
\If{no supporting segment after a completed search}
  \State \Return scoped not-found (searched scope, blind spots)
\Else
  \State \Return answer with provenance (top admitted segments)
\EndIf
\end{algorithmic}
\end{algorithm}
Exhausting a policy's call budget is intended to produce a scope and blind-spot
report. The endpoint enforces timeouts and turn limits; the per-policy budgets
are workflow instructions, not server quotas.

\subsection{Tier policies and the evidence gate}
\label{sec:trace}
\label{sec:gate}
\label{sec:gate-primitive}
\label{sec:twostage}

The chat endpoint supplies two built-in policies. Lane A uses a wide dense
search (top 30, cosine threshold 0.35) within a pinned collection and may
escalate once to the catalogue ladder when no verified evidence survives;
because the collection is pinned, the catalogue tier is bypassed rather than
traversed. Lane B follows the catalogue tier and then the shelf tier, and
uses a deterministic scanner for literal or enumeration questions and the
judge for semantic questions. Endpoint-level timeouts and turn limits bound
execution; the lane call budgets and six-reference batching below are workflow
instructions, not independent server-enforced quotas. Exhaustion is intended
to produce a scope and blind-spot report. Lane C is a composed configuration
used in evaluation, not a third built-in server policy: it appends a hybrid
directive to the chat prompt that can call vector and catalogue tools, merge up
to twelve candidates, judge them in batches of six, and return up to five
survivors with bounded evidence. Table~\ref{tab:lanes} summarises the
tier-specific policies; Figure~\ref{fig:workflow} shows their place in the
full document-to-answer workflow.

The gate is the local Laya evidence classifier \cite{laya}, exposed as
\path{kb_laya_judge}. In reference mode, the server fetches each document,
judges its segments and returns provenance with a bounded evidence pack. The
verdict controls admission, not truth; the verifier remains subject to
calibration and transport limitations.

\begin{table*}[t]
\centering
\small
\setlength{\tabcolsep}{5pt}
\caption{Access policies over the tiered path. Policies differ in which tier
performs the narrowing, not in the ranking model at the bottom of the path.}
\label{tab:lanes}
\begin{tabular}{@{}>{\raggedright\arraybackslash}p{0.11\textwidth}>{\raggedright\arraybackslash}p{0.19\textwidth}>{\raggedright\arraybackslash}p{0.24\textwidth}>{\raggedright\arraybackslash}p{0.16\textwidth}>{\raggedright\arraybackslash}p{0.20\textwidth}@{}}
\toprule
Policy & Tiers traversed & Retrieval affordance & Budget type & Intended terminal report\\
\midrule
Search & pinned collection, segment & Dense search within one collection & Endpoint turn limit & Answer or report not-found\\
Librarian & catalogue, shelf, document, segment & Catalogue navigation and query-shaped scan/judging & Workflow call budget & Scope and blind spots\\
Hybrid & collection, segment & Agent-composed vector and catalogue calls & Evaluation-specific candidate cap & Survivors for answer workflow\\
\bottomrule
\end{tabular}
\end{table*}

The two-stage tool builds a candidate pool with lexical search and optional
graph expansion, then applies dense ranking. It does not call the
evidence gate; a ranked chunk is not admissible evidence until verified, and a
similarity threshold is not an abstention mechanism. The 13 October catalogue
incident showed that oversized tool returns can conceal documents before
reasoning begins. The tested repair adds a server-side character budget,
grouped views and pagination; its behaviour is reported in the Evaluation
section.

\subsection{Honest abstention}
A candidate list may contain high-similarity but unsupported documents. The
answer workflow instructs the agent to report not-found when its completed
evidence review finds no support, including the searched scope and known blind
spots, rather than falling back to parametric knowledge. This behaviour was
observed in the reported trials; it is not a server-side semantic-classifier
guarantee. Timeouts, unavailable tools and incomplete searches are execution
failures, not evidence of semantic absence (Section~\ref{sec:eval}).
